% Generated draft: requires mathematical and source review.
\documentclass[11pt]{article}
\usepackage[margin=25mm]{geometry}
\usepackage{fontspec}
\setmainfont{Latin Modern Roman}
\usepackage{amsmath,amssymb,mathtools,unicode-math}
\setmathfont{Latin Modern Math}
\usepackage{enumitem,needspace,xurl,hyperref}
\setlength{\parindent}{0pt}
\setlength{\parskip}{5pt}
\setlength{\emergencystretch}{3em}
\newcommand{\N}{\mathbb{N}}
\newcommand{\Z}{\mathbb{Z}}
\newcommand{\Q}{\mathbb{Q}}
\newcommand{\R}{\mathbb{R}}
\newcommand{\C}{\mathbb{C}}
\newcommand{\F}{\mathbb{F}}
\newcommand{\A}{\mathbb{A}}
\newcommand{\PP}{\mathbb{P}}
\newcommand{\E}{\mathbb{E}}
\newcommand{\eps}{\varepsilon}
\newcommand{\dd}{\,\mathrm{d}}
\begin{document}
\section*{UnsolvedMath ID: 30004535}
\textbf{Source problem code:} catalogue entry
\subsection{Definitions and mathematical statement}
The source asks about a sequence of ideals \(I_n\subseteq R_n\), where \(R_n\) are polynomial rings and the sequence is described as a \(\operatorname{Sym}\)-invariant chain. The record does not specify the coefficient field, the variables of \(R_n\), the transition maps between rings, the symmetric-group actions, or the precise compatibility conditions defining such a chain. Thus the class of chains quantified over below remains dependent on these missing definitions.

For a nonzero module \(M\) over a ring \(R\), its projective dimension \(\operatorname{pd}_R(M)\) is the least length of a projective resolution, or \(+\infty\) if no finite projective resolution exists. To explain the regularity notation conditionally, if \(R_n\) is a standard graded polynomial ring over a field \(k\) and \(I_n\) is a nonzero homogeneous ideal, its Castelnuovo--Mumford regularity as an \(R_n\)-module is
\[
\operatorname{reg}_{R_n}(I_n)
=\max\{j-i:\operatorname{Tor}^{R_n}_i(I_n,k)_j\ne 0\}.
\]
Here the subscript \(j\) denotes internal degree. The source writes this invariant as \(\operatorname{reg}(I_n)\). The standard grading and homogeneity assumptions needed for this interpretation are not explicitly supplied; they require confirmation rather than silent addition to the problem.

\textbf{Eventual-linearity question, subject to these unresolved specifications.} For every \(\operatorname{Sym}\)-invariant chain \((I_n\subseteq R_n)\) in the intended class, do there exist integers \(A,B,C,D\) and an index \(N\) such that, for every index \(n\ge N\),
\[
\operatorname{pd}_{R_n}(R_n/I_n)=An+B,
\qquad
\operatorname{reg}_{R_n}(I_n)=Cn+D?
\]
The integers and the threshold may depend on the chain. Here ``linear'' means affine linear, as specified by the displayed formulas. The source concerns the projective dimension of the quotient and the regularity of the ideal. It supplies no exclusions or conventions for \(I_n=0\) or \(I_n=R_n\); these cases must be resolved before the question is a fully specified assertion with integer-valued invariants.
\subsection{Short English statement}
The question asks whether every intended symmetric chain of polynomial ideals eventually has affine-linear projective dimension for its quotients and affine-linear regularity for its ideals. The ambient rings, symmetry and chain conditions, grading assumptions, and treatment of zero modules need clarification. The supplied status is partially solved; this is a dated curation label, not an independently verified mathematical conclusion.
\subsection{Sources}
\begin{itemize}
\item Local UnsolvedMath registry: supplied complete record for permanent ID \(30004535\), source problem code \texttt{catalogue entry}, titled ``Eventual Linearity of Homological Invariants in Symmetric Ideal Chains,'' in the Oberwolfach Reports Open Problems collection.
\item Supplied report citation: \emph{Lefschetz Properties in Algebra, Geometry and Combinatorics} (2021), Oberwolfach Reports, pp.\ 1539--1600, DOI \texttt{10.4171/owr/2020/31}. Supplied link: \url{https://doi.org/10.4171/owr/2020/31}. The registry records proposed year 2020 and names no individual proposer. It describes Oberwolfach Reports as documenting mathematical workshops and their open-problem discussions; no problem-specific page is supplied.
\item Supplied literature reference labelled ``Le--Nguyen'': \url{https://arxiv.org/abs/2206.15141}.
\item Supplied literature reference labelled ``Regularity conjecture'': \url{https://arxiv.org/abs/1806.00457}.
\item Supplied literature reference labelled ``Projective-dimension conjecture'': \url{https://arxiv.org/abs/1809.06877}.
\item Supplied literature assessment, checked at the exact recorded date \texttt{2026-08-22}: eventual linearity of projective dimension and regularity is reported as proved for several monomial, saturated, and related special classes, while the full conjecture for \(\operatorname{Sym}\)-invariant chains is reported as remaining open. This is supplied dated curation derived from the OWR material, not an independently verified assessment; the time-sensitive status and the precise scope of the cited results require independent verification.
\item The registry labels the statement \texttt{corrected\_verified} and reports that the official report was checked and the symmetric-chain setting and both eventual-linearity formulas were restored in publication notation. This is a supplied source-verification claim, not a verification performed for this draft.
\end{itemize}
\end{document}
